\section{The mechanism of the counterexample}
\label{sec:structure}
The construction uses the coordinates $r(a)$ and $ma$ to express the
symmetric part of the graph pairing:
\[
\pair{La-Lb}{a-b}=|r(a)-r(b)|^2-\|ma-mb\|_2^2.
\]
A Lipschitz curve in these coordinates makes the graph monotone.
Maximality requires the additional information supplied by the complete
affine fibres. Testing every kernel direction forces an arbitrary polar
query to have the scalar form in Eq.~\eqref{eq:abs-polar-square}.
At an exterior parameter this form returns an actual graph point.
Over the middle interval, the two endpoint comparisons force the
completed profile $(\tau,\eta)$.

Theorem~\ref{thm:endpoint-criterion} gives the exact distinction:
the completed profile contributes a polar point precisely when
some $p\in X^*$ satisfies both $r(p)=\tau$ and $Mp=(\tau,\eta)$.
For the constructed operator, $mp\in\ell^1$ and
$\eta=(1/(4n))_{n\ge1}\notin\ell^1$. Thus the comparison geometry
has a middle segment with no continuous-dual labels realizing it.
The proof takes limits in the comparison space; it never enlarges
the dual space of $c_0$.

This failure of realization can also be seen along actual graph rows.
Set $t_n^\pm=\pm(1+1/n)$, $a_n^\pm=a^{t_n^\pm}$, and
$x_n^\pm=La_n^\pm$. For $t\in J$ and $\epsilon=\sgn t$,
direct substitution in Eqs.~\eqref{eq:D4}--\eqref{eq:D9} gives
\[
(La^t)_{0,1}=-2t-2\epsilon,\qquad
(La^t)_{0,2}=-3t-2\epsilon,\qquad
(La^t)_{0,3}=-t-\epsilon,
\]
with the remaining block-zero coordinates zero. In block $b\ge1$,
the only nonzero coordinate is $-\omega_{|t|-1}(b)$ at index $1$.
Consequently, $x_n^\pm$ converge in the supremum norm to
$x_\pm\in c_0(I)$ with block-zero coordinates $\mp(4,5,2)$ and
positive-block first coordinates $-1/(4b)$. Indeed, convergence of
the tails in $\ell^2$ implies convergence in the supremum norm.
The limits satisfy $\pair{x_\pm}{g}=\pm1$ and lie outside $\dom A$.
On the other hand,
\[
\|a_n^\pm\|_1=2(1+1/n)+1+\|\omega_{1/n}\|_1\longrightarrow\infty.
\]
For each fixed $N$, the sum of the first $N$ tail coordinates tends to
$\sum_{b=1}^N1/(4b)$; these partial sums are unbounded. This proves
the divergence without asserting convergence of the labels in the dual.

The sum obstruction uses a different consequence of the pairing identity.
Every row of the first graph satisfies
\[
\pair{La}{a}=r(a)^2-1-W(|r(a)|-1)>-\sigma,\qquad
\|La\|_\infty>1/2.
\]
These inequalities prove $\mathcal V(\gra A)\le2\sigma$ directly.
The second factor adds exactly $r(a)^2$, making the pairing positive
on the entire sum graph. Since zero remains outside the domain,
$(0,0)$ is a missing polar point of the sum. The finite radial bound
is thus an output of the construction, not an assumed universal
property or an unproved consequence of the desired counterexample.
